\section{Conclusion}
\label{sec:summary}

Software testing research has long relied on structural metrics as proxies for test suite quality, partly because reliable approaches do not exist for measuring whether a test suite validates a method's documented behaviour.
In this paper, we introduce the \toolname{} proof-of-concept that extracts method-level expected behaviours from documentation and source code, maps them to test cases, and identifies which documented behaviours go untested. 
Across 8,922 methods in ten well-tested Java libraries, \toolname{} surfaces 20,729 behaviours at $93.1\%\pm2.7\%$ precision. At least 17.5\% of those detected behaviours are untested in developer-written suites. Test suites generated by ASTER / EvoSuite leave 27.1\% / 20.6\% of detected behaviours untested, both exceeding developer-written behavioural gaps on the same projects. 
Qualitative inspection of the missed behaviours (exception handling, side effects, null input handling) confirms that none are inherently untestable.
Behavioural gaps provide useful independent insight from structural metrics: Among methods with perfect line coverage in our dataset, 38.2\% still contain at least one behavioural gap; among methods with perfect mutation kill scores, 29.6\% do. Spearman correlations of 0.112 (line) and 0.377 (mutation) confirm that neither metric reliably predicts behavioural gaps. 
Ultimately, 
coverage reports and kill scores tell developers how thoroughly the implementation has been exercised, but behavioural gap analysis can provide novel test adequacy insight into whether code was validated consistently with its expected behaviour.

